\documentclass[10pt,a4paper]{amsart}
\usepackage[margin=19mm]{geometry}
\usepackage{amsmath,amssymb,mathtools}
\usepackage[scr=rsfs,cal=euler]{mathalfa}
\usepackage{xfrac}
\usepackage[unicode,hidelinks,bookmarks=false]{hyperref}
\newcommand{\ie}{i.e.\ }
\newcommand{\cf}{cf.\ }
\newcommand{\resp}{respectively\ }
\newcommand{\Subsection}{\subsection}
\providecommand{\nobreakdash}{\nobreak\hbox{-}}
\setlength{\emergencystretch}{2em}
\multlinegap0pt
\usepackage{mathtools}
\usepackage{amssymb}
\newtheorem{lemma}{Lemma}
\newcommand{\V}{V}
\title{The Hodge conjecture for Fermat fourfolds of odd degree at most 199}
\author{Rifat Jumagulov}
\date{Source: arXiv:2608.18134v1 (CC BY 4.0)}
\begin{document}
\maketitle

\noindent\emph{Source and licence.} The Hodge conjecture for Fermat fourfolds of odd degree at most 199. Version arXiv:2608.18134v1 (CC BY 4.0), distributed by arXiv under the Creative Commons Attribution 4.0 International licence, http://creativecommons.org/licenses/by/4.0/, as recorded in the arXiv licence metadata for this exact version and shown on the abstract page https://arxiv.org/abs/2608.18134v1. Full licence notice: \texttt{licenses/arxiv-2608.18134-CC-BY-4.0.txt}.\par\smallskip

\noindent\textit{Editorial scope.} Counted unit: the article's lemma lem:descent (source0.tex lines 550--554). The article's complete original proof of it is delivered with it (source0.tex lines 555--563, one \texttt{proof} environment of 9 lines). No mathematics is written, repaired or re-derived: the statement and the proof are the article's own lines, in the article's own notation, and no retained line is substituted or re-flowed.  No further source-local context is needed: every cross-reference of every retained line resolves inside the two delivered spans.  Nothing was omitted from any retained span, so the package carries no omission record.  No retained line contains a citation, so the wrapper carries no bibliography and no bibliography entry.  No retained line contains a figure environment, an image-inclusion command or drawing code, so no image is delivered and the package declares no asset dependency.  Editorial formatting only, in addition: an AMS \texttt{amsart} preamble that restates the article's own declarations and macros used by the retained lines, namely the environments \texttt{lemma} on separate counters, together with the standard packages those lines need.  The retained environments print in this document as Lemma 1 in order of appearance, whereas the article prints the same items with the numbering of its own full document.  The complete native source of the article as distributed in the arXiv e-print for this exact version is bundled unchanged as \texttt{sources/207971/source0.tex}.
\begin{lemma}[descent]\label{lem:descent}
Let $m=e\,m'$ and $\beta$ a character of level $m'$ (equivalently: $e\beta$ has content $e$
and all entries nonzero mod $m$). Then
$\mathrm{claim}_m(e\beta)\Rightarrow\mathrm{claim}_{m'}(\beta)$.
\end{lemma}

\begin{proof}
$\pi\colon X^n_m\to X^n_{m'}$, $(x_i)\mapsto(x_i^e)$ is a finite surjective morphism
($\sum x_i^m=\sum(x_i^e)^{m'}$ identically) intertwining the group actions along the $e$-th-power
surjection $G_m\twoheadrightarrow G_{m'}$; hence $\pi_*\V(e\beta)\subseteq\V(\beta)$ (again: the entries of $e\beta$ are nonzero, so $\V(e\beta)$ exists). Since
$\pi_*\pi^*=(\deg\pi)\,\mathrm{id}\neq0$ and $\pi^*\V(\beta)=\V(e\beta)$ with both eigenlines
$1$-dimensional, $\pi_*\colon\V(e\beta)\xrightarrow{\ \sim\ }\V(\beta)$. Proper pushforward carries
algebraic cycle classes to algebraic cycle classes ($\mathrm{cl}\circ\pi_*=\pi_*\circ\mathrm{cl}$),
and $\pi$ is finite of relative dimension zero, so there is no degree shift or Tate twist.
\end{proof}

\end{document}
